\documentclass[10pt,aspectratio=169]{beamer}

% ------------------------------------------------------------------
% Presentació adaptada: Probabilitat 1 amb orientació computacional
% Grau en Enginyeria Informàtica - Estadística
% ------------------------------------------------------------------

\usepackage[utf8]{inputenc}
\usepackage[T1]{fontenc}
\usepackage[catalan]{babel}
\usepackage{amsmath,amssymb}
\usepackage{booktabs}
\usepackage{array}
\usepackage{tikz}
\usepackage{pgfplots}
\usepackage{xcolor}
\usepackage{listings}
\usetikzlibrary{positioning,arrows.meta,calc,shapes.geometric}
\pgfplotsset{compat=1.18}

% ------------------------------------------------------------------
% Estil visual coherent amb les presentacions anteriors
% ------------------------------------------------------------------
\definecolor{UABGreen}{RGB}{0,103,71}
\definecolor{UABLight}{RGB}{224,241,236}
\definecolor{DeepBlue}{RGB}{20,65,110}
\definecolor{AccentOrange}{RGB}{230,126,34}
\definecolor{SoftGray}{RGB}{245,247,250}
\definecolor{CodeBack}{RGB}{248,248,248}
\definecolor{CodeFrame}{RGB}{210,215,220}

\usetheme{Madrid}
\usecolortheme[named=UABGreen]{structure}
\setbeamercolor{title}{fg=white,bg=UABGreen}
\setbeamercolor{frametitle}{fg=white,bg=UABGreen}
\setbeamercolor{block title}{fg=white,bg=DeepBlue}
\setbeamercolor{block body}{fg=black,bg=SoftGray}
\setbeamercolor{alerted text}{fg=AccentOrange}
\setbeamertemplate{navigation symbols}{}
\setbeamertemplate{itemize item}{\color{UABGreen}$\blacktriangleright$}
\setbeamertemplate{itemize subitem}{\color{AccentOrange}$\bullet$}
\setbeamertemplate{enumerate item}{\color{UABGreen}\insertenumlabel.}

\setbeamertemplate{footline}{%
  \leavevmode%
  \hbox{%
  \begin{beamercolorbox}[wd=.72\paperwidth,ht=2.4ex,dp=1ex,leftskip=1.4em]{author in head/foot}%
    Estadística -- Grau en Enginyeria Informàtica
  \end{beamercolorbox}%
  \begin{beamercolorbox}[wd=.28\paperwidth,ht=2.4ex,dp=1ex,center]{date in head/foot}%
    \insertframenumber{} / \inserttotalframenumber
  \end{beamercolorbox}}%
  \vskip0pt%
}

\lstdefinelanguage{R}{%
  morekeywords={if,else,repeat,while,function,for,in,next,break,TRUE,FALSE,NULL,NA,NaN,Inf,
    library,mean,median,sd,var,summary,table,prop.table,xtabs,addmargins,round,
    ggplot,aes,geom_point,geom_line,geom_col,geom_tile,geom_smooth,facet_wrap,
    labs,theme_minimal,summarise,group_by,mutate,read_csv,filter,arrange,set.seed,
    rbinom,runif,sample,replicate,cumsum,cummean,tibble,rep,seq,case_when,
    future_map_dbl,plan,multisession,furrr,pmap_dbl,rowwise},
  sensitive=true,
  morecomment=[l]\#,
  morestring=[b]",
  morestring=[b]'
}
\lstset{%
  language=R,
  basicstyle=\ttfamily\scriptsize,
  keywordstyle=\color{DeepBlue}\bfseries,
  commentstyle=\color{gray},
  stringstyle=\color{AccentOrange},
  backgroundcolor=\color{CodeBack},
  frame=single,
  rulecolor=\color{CodeFrame},
  breaklines=true,
  columns=fullflexible,
  keepspaces=true,
  showstringspaces=false,
  literate={à}{{\`a}}1 {è}{{\`e}}1 {é}{{\'e}}1 {í}{{\'i}}1 {ï}{{\"i}}1 {ò}{{\`o}}1 {ó}{{\'o}}1 {ú}{{\'u}}1 {ü}{{\"u}}1 {ç}{{\c{c}}}1 {À}{{\`A}}1 {È}{{\`E}}1 {É}{{\'E}}1 {Í}{{\'I}}1 {Ò}{{\`O}}1 {Ó}{{\'O}}1 {Ú}{{\'U}}1
}

\newcommand{\R}{\textsf{R}}
\newcommand{\RStudio}{\textsf{RStudio/Posit}}
\newcommand{\GitHub}{\textsf{GitHub}}
\newcommand{\important}[1]{\textcolor{UABGreen}{\textbf{#1}}}
\newcommand{\exemple}[1]{\textcolor{AccentOrange}{\textbf{#1}}}
\newcommand{\Prob}{\mathbb{P}}

% ------------------------------------------------------------------
% Dades de la presentació
% ------------------------------------------------------------------
\title[Probabilitat computacional]{Probabilitat I}
\author{Estadística -- Grau en Enginyeria Informàtica}
\date{Curs 2026--27}

\begin{document}

% ------------------------------------------------------------------
\begin{frame}[plain]
  \titlepage
  \vspace{-0.5cm}
  \begin{center}
  \small
  \textcolor{DeepBlue}{La probabilitat com a llenguatge per raonar amb dades, algoritmes i sistemes incerts}
  \end{center}
\end{frame}

% ------------------------------------------------------------------
\begin{frame}{Per què probabilitat en Enginyeria Informàtica?}
\begin{columns}[T]
\column{0.56\textwidth}
\begin{itemize}
  \item Molts sistemes informàtics no són deterministes a escala observable.
  \item Treballem amb \important{errors, fallades, latències, mostres, classificadors, soroll i decisions sota incertesa}.
  \item La probabilitat permet modelitzar:
  \begin{itemize}
    \item pèrdua de paquets en una xarxa,
    \item falsos positius en un detector,
    \item col·lisions en hashing,
    \item experiments A/B,
    \item filtres de correu brossa,
    \item fiabilitat de sistemes distribuïts.
  \end{itemize}
\end{itemize}
\column{0.38\textwidth}
\begin{block}{Idea clau}
No n'hi ha prou amb dir ``ha passat'' o ``no ha passat''. Volem quantificar \important{com de probable} és i com canvia aquesta probabilitat quan tenim nova informació.
\end{block}
\vspace{0.2cm}
\begin{tikzpicture}
\node[draw=UABGreen,rounded corners,fill=UABLight,text width=4.4cm,align=center,inner sep=8pt] {
Dades $+$ model $+$ simulació $\rightarrow$ decisió informada
};
\end{tikzpicture}
\end{columns}
\end{frame}

% ------------------------------------------------------------------
\begin{frame}{Objectius del tema}
Al final d'aquest tema hauríeu de poder:
\begin{enumerate}
  \item Interpretar la probabilitat com a freqüència relativa a llarg termini i com a model.
  \item Definir espai mostral, esdeveniments, complementaris, unions i interseccions.
  \item Calcular probabilitats a partir de taules de freqüències.
  \item Entendre i aplicar la probabilitat condicionada.
  \item Identificar esdeveniments independents.
  \item Aplicar la fórmula de probabilitat total i el teorema de Bayes.
  \item Verificar resultats amb simulació en \R.
  \item Connectar aquests conceptes amb classificació, detecció d'errors i monitoratge de sistemes.
\end{enumerate}
\end{frame}

% ------------------------------------------------------------------
\begin{frame}{Mapa del tema}
\begin{center}
\begin{tikzpicture}[node distance=1.25cm, every node/.style={font=\small}]
\node[draw,rounded corners,fill=UABLight,text width=3.0cm,align=center,minimum height=0.75cm] (freq) {Freqüències relatives};
\node[draw,rounded corners,fill=SoftGray,below of=freq,yshift=-0.1cm,text width=3.0cm,align=center,minimum height=0.75cm] (prob) {Probabilitat formal};
\node[draw,rounded corners,fill=SoftGray,below of=prob,yshift=-0.1cm,text width=3.0cm,align=center,minimum height=0.75cm] (cond) {Probabilitat condicionada};
\node[draw,rounded corners,fill=SoftGray,right of=cond,xshift=3.2cm,text width=3.0cm,align=center,minimum height=0.75cm] (ind) {Independència};
\node[draw,rounded corners,fill=SoftGray,below of=cond,yshift=-0.1cm,text width=3.0cm,align=center,minimum height=0.75cm] (total) {Probabilitat total};
\node[draw,rounded corners,fill=UABLight,right of=total,xshift=3.2cm,text width=3.0cm,align=center,minimum height=0.75cm] (bayes) {Bayes};
\node[draw,rounded corners,fill=AccentOrange!15,below of=total,yshift=-0.1cm,text width=6.8cm,align=center,minimum height=0.75cm] (sim) {Simulació amb R i interpretació computacional};
\draw[->,thick,UABGreen] (freq) -- (prob);
\draw[->,thick,UABGreen] (prob) -- (cond);
\draw[->,thick,UABGreen] (cond) -- (total);
\draw[->,thick,UABGreen] (total) -- (bayes);
\draw[->,thick,AccentOrange] (cond) -- (ind);
\draw[->,thick,AccentOrange] (bayes) -- (sim);
\draw[->,thick,AccentOrange] (total) -- (sim);
\end{tikzpicture}
\end{center}
\vspace{0.1cm}
\begin{alertblock}{Pregunta transversal}
Com actualitzem una probabilitat quan observem una nova evidència?
\end{alertblock}
\end{frame}

% ------------------------------------------------------------------
\section{Probabilitat i simulació}

\begin{frame}{Un problema històric: de Méré i Fermat}
\begin{columns}[T]
\column{0.52\textwidth}
\begin{block}{Joc 1}
Llancem un dau 4 vegades. Guanyem si surt almenys un 6.
\[
\Prob(\text{guanyar})=1-\left(\frac{5}{6}\right)^4=0.5177.
\]
\end{block}
\begin{block}{Joc 2}
Llancem dos daus 24 vegades. Guanyem si surt almenys un doble 6.
\[
\Prob(\text{guanyar})=1-\left(\frac{35}{36}\right)^{24}=0.4914.
\]
\end{block}
\column{0.42\textwidth}
\begin{alertblock}{Lliçó}
Una intuïció aparentment raonable pot fallar. La probabilitat ens dona una manera formal de comprovar-la.
\end{alertblock}
\vspace{0.2cm}
\begin{tikzpicture}
\node[draw=UABGreen,rounded corners,fill=UABLight,text width=4.7cm,align=center,inner sep=8pt] {
Càlcul exacte $+$ simulació $\Rightarrow$ comprensió del fenomen
};
\end{tikzpicture}
\end{columns}
\end{frame}

% ------------------------------------------------------------------
\begin{frame}[fragile]{El mateix problema amb R}
\begin{lstlisting}
# Probabilitats exactes
p_joc1 <- 1 - (5/6)^4
p_joc2 <- 1 - (35/36)^24

p_joc1
p_joc2

# Simulació del segon joc
set.seed(123)
B <- 100000
guanya <- replicate(B, {
  daus1 <- sample(1:6, size = 24, replace = TRUE)
  daus2 <- sample(1:6, size = 24, replace = TRUE)
  any(daus1 == 6 & daus2 == 6)
})
mean(guanya)
\end{lstlisting}
\vspace{0.2cm}
\begin{block}{Missatge computacional}
La simulació no substitueix la teoria, però ajuda a validar càlculs, detectar errors d'intuïció i entendre la variabilitat.
\end{block}
\end{frame}

% ------------------------------------------------------------------
\begin{frame}{Freqüències relatives i probabilitats}
\begin{columns}[T]
\column{0.52\textwidth}
Per esdeveniments aleatoris independents repetits moltes vegades:
\[
\text{freqüència relativa} = \frac{\text{nombre d'èxits}}{\text{nombre de repeticions}}
\]
\vspace{-0.1cm}
\[
\text{freqüència relativa} \approx \Prob(\text{èxit}).
\]
\begin{itemize}
  \item Amb poques repeticions pot variar força.
  \item Amb moltes repeticions tendeix a estabilitzar-se.
  \item Aquesta és la base de la simulació Monte Carlo.
\end{itemize}
\column{0.43\textwidth}
\begin{tikzpicture}
\begin{axis}[
  width=6.0cm,height=4.2cm,
  xlabel={Nombre de simulacions}, ylabel={Freq. relativa},
  ymin=0.35,ymax=0.65,xmin=0,xmax=1000,
  axis lines=left, grid=major, grid style={gray!20},
  tick label style={font=\scriptsize}, label style={font=\scriptsize}]
\addplot[DeepBlue, thick] coordinates {
(10,0.600)(25,0.520)(50,0.540)(100,0.500)(150,0.507)(200,0.515)(300,0.503)(400,0.510)(600,0.505)(800,0.496)(1000,0.492)};
\addplot[AccentOrange, dashed, thick] coordinates {(0,0.4914)(1000,0.4914)};
\end{axis}
\end{tikzpicture}
\end{columns}
\end{frame}

% ------------------------------------------------------------------
\begin{frame}[fragile]{Un patró de simulació reutilitzable}
\begin{columns}[T]
\column{0.58\textwidth}
\begin{lstlisting}
set.seed(2026)
B <- 5000

simula_joc <- function() {
  daus1 <- sample(1:6, 24, replace = TRUE)
  daus2 <- sample(1:6, 24, replace = TRUE)
  any(daus1 == 6 & daus2 == 6)
}

resultats <- replicate(B, simula_joc())
freq_acumulada <- cumsum(resultats) / seq_len(B)
\end{lstlisting}
\column{0.38\textwidth}
\begin{block}{Bones pràctiques}
\begin{itemize}
  \item Fixar una llavor amb \texttt{set.seed()}.
  \item Escriure funcions petites.
  \item Separar simulació, càlcul i visualització.
  \item Guardar el codi en un projecte de \RStudio{} i \GitHub{}.
\end{itemize}
\end{block}
\end{columns}
\end{frame}

% ------------------------------------------------------------------
\section{Probabilitat a partir de dades}

\begin{frame}{Esdeveniments a partir d'una taula de dades}
Suposem un registre de 120 peticions a un sistema web:
\begin{center}
\small
\begin{tabular}{lccc}
\toprule
 & Error & No error & Total \\
\midrule
Mòbil & 25 & 50 & 75 \\
Escriptori & 20 & 25 & 45 \\
\midrule
Total & 45 & 75 & 120 \\
\bottomrule
\end{tabular}
\end{center}
\vspace{0.2cm}
\begin{columns}[T]
\column{0.48\textwidth}
\begin{itemize}
  \item $E$: la petició acaba amb error.
  \item $M$: la petició ve de mòbil.
  \item $D$: la petició ve d'escriptori.
\end{itemize}
\column{0.48\textwidth}
\begin{align*}
\Prob(E) &= 45/120 = 0.375,\\
\Prob(M) &= 75/120 = 0.625,\\
\Prob(D \cap E) &= 20/120 = 0.1667.
\end{align*}
\end{columns}
\end{frame}

% ------------------------------------------------------------------
\begin{frame}{Unions, interseccions i complementaris}
\begin{columns}[T]
\column{0.52\textwidth}
\begin{block}{Regla de la unió}
\[
\Prob(A\cup B)=\Prob(A)+\Prob(B)-\Prob(A\cap B).
\]
\end{block}
\begin{block}{Complementari}
\[
\Prob(\overline{A})=1-\Prob(A).
\]
\end{block}
\begin{exampleblock}{Exemple}
\[
\Prob(D\cup E)=0.375+0.375-0.1667=0.5833.
\]
\end{exampleblock}
\column{0.42\textwidth}
\begin{center}
\begin{tikzpicture}[scale=0.9]
\fill[UABLight] (-1,0) circle (1.15);
\fill[AccentOrange!25] (1,0) circle (1.15);
\draw[thick,DeepBlue] (-1,0) circle (1.15);
\draw[thick,UABGreen] (1,0) circle (1.15);
\node at (-1.35,0.9) {$A$};
\node at (1.35,0.9) {$B$};
\node at (0,-0.05) {$A\cap B$};
\node[align=center] at (0,-1.8) {La intersecció no s'ha de comptar dues vegades};
\end{tikzpicture}
\end{center}
\end{columns}
\end{frame}

% ------------------------------------------------------------------
\begin{frame}[fragile]{Taules de probabilitats amb R}
\begin{lstlisting}
peticions <- tibble::tribble(
  ~dispositiu,   ~estat,      ~n,
  "mobil",      "error",      25,
  "mobil",      "no_error",   50,
  "escriptori", "error",      20,
  "escriptori", "no_error",   25
)

# Taula de freqüències absolutes
xt <- xtabs(n ~ dispositiu + estat, data = peticions)
addmargins(xt)

# Taula de probabilitats relatives
round(prop.table(xt), 4)

# Probabilitats condicionades per dispositiu
round(prop.table(xt, margin = 1), 4)
\end{lstlisting}
\vspace{0.2cm}
\begin{block}{Connexió amb descriptiva bivariant}
Les taules de contingència també són una porta d'entrada natural a la probabilitat.
\end{block}
\end{frame}

% ------------------------------------------------------------------
\section{Definició formal}

\begin{frame}{Definició formal de probabilitat}
\begin{columns}[T]
\column{0.52\textwidth}
\begin{itemize}
  \item \important{Espai mostral} $\Omega$: conjunt de tots els resultats possibles.
  \item \important{Esdeveniments}: subconjunts de $\Omega$.
  \item Una probabilitat és una funció
  \[
  \Prob: \mathcal{E} \longrightarrow [0,1]
  \]
  que assigna un nombre a cada esdeveniment.
\end{itemize}
\column{0.42\textwidth}
\begin{block}{Axiomes bàsics}
\begin{enumerate}
  \item $0\leq \Prob(A)\leq 1$.
  \item $\Prob(\Omega)=1$ i $\Prob(\emptyset)=0$.
  \item Si $A$ i $B$ són disjunts:
  \[
  \Prob(A\cup B)=\Prob(A)+\Prob(B).
  \]
\end{enumerate}
\end{block}
\end{columns}
\vspace{0.1cm}
\begin{alertblock}{Interpretació computacional}
Un esdeveniment pot ser: ``temps de resposta $>$ 1 segon'', ``paquet perdut'', ``classificador diu spam'', ``test automàtic falla''.
\end{alertblock}
\end{frame}

% ------------------------------------------------------------------
\begin{frame}{Espai mostral en un sistema informàtic}
\begin{center}
\begin{tikzpicture}[node distance=1cm, every node/.style={font=\small}]
\node[draw,rounded corners,fill=UABLight,text width=3.3cm,align=center,minimum height=0.8cm] (omega) {$\Omega$\\Totes les peticions};
\node[draw,rounded corners,fill=SoftGray,below left=of omega,xshift=-0.5cm,text width=3.0cm,align=center,minimum height=0.8cm] (ok) {Peticions OK};
\node[draw,rounded corners,fill=SoftGray,below right=of omega,xshift=0.5cm,text width=3.0cm,align=center,minimum height=0.8cm] (err) {Peticions amb error};
\node[draw,rounded corners,fill=AccentOrange!15,below=of err,text width=3.0cm,align=center,minimum height=0.8cm] (timeout) {Timeout};
\node[draw,rounded corners,fill=AccentOrange!15,below=of ok,text width=3.0cm,align=center,minimum height=0.8cm] (fast) {Latència $<100$ ms};
\draw[->,thick,UABGreen] (omega) -- (ok);
\draw[->,thick,UABGreen] (omega) -- (err);
\draw[->,thick,AccentOrange] (err) -- (timeout);
\draw[->,thick,AccentOrange] (ok) -- (fast);
\end{tikzpicture}
\end{center}
\vspace{0.2cm}
\begin{block}{Traducció estadística}
Cada filtre sobre les dades defineix un esdeveniment. Les probabilitats poden estimar-se com a freqüències relatives o modelitzar-se teòricament.
\end{block}
\end{frame}

% ------------------------------------------------------------------
\section{Probabilitat condicionada i independència}

\begin{frame}{Probabilitat condicionada}
\begin{block}{Definició}
La probabilitat que passi $A$ sabent que ha passat $B$ és:
\[
\Prob(A\mid B)=\frac{\Prob(A\cap B)}{\Prob(B)}, \qquad \Prob(B)>0.
\]
\end{block}
\vspace{0.2cm}
\begin{columns}[T]
\column{0.48\textwidth}
\begin{exampleblock}{Exemple}
Probabilitat que una petició tingui error sabent que ve d'escriptori:
\[
\Prob(E\mid D)=\frac{\Prob(E\cap D)}{\Prob(D)}
=\frac{20/120}{45/120}=0.4444.
\]
\end{exampleblock}
\column{0.48\textwidth}
\begin{alertblock}{Interpretació}
Condicionar vol dir canviar la població de referència. Ja no mirem totes les peticions, només les que compleixen $D$.
\end{alertblock}
\end{columns}
\end{frame}

% ------------------------------------------------------------------
\begin{frame}{Independència}
\begin{block}{Definició}
Dos esdeveniments $A$ i $B$ són independents si:
\[
\Prob(A\cap B)=\Prob(A)\Prob(B).
\]
Equivalentment, si $\Prob(B)>0$:
\[
\Prob(A\mid B)=\Prob(A).
\]
\end{block}
\vspace{0.2cm}
\begin{columns}[T]
\column{0.48\textwidth}
\begin{itemize}
  \item Si saber $B$ no canvia la probabilitat de $A$, no hi ha informació addicional.
  \item En dades reals, la independència exacta és rara.
\end{itemize}
\column{0.48\textwidth}
\begin{alertblock}{En informàtica}
No assumiu independència sense pensar-hi: errors de xarxa, caigudes de serveis i temps de resposta poden estar correlacionats.
\end{alertblock}
\end{columns}
\end{frame}

% ------------------------------------------------------------------
\begin{frame}[fragile]{Comprovar independència amb dades}
\begin{lstlisting}
# Probabilitat d'error global
p_error <- sum(xt[, "error"]) / sum(xt)

# Probabilitat d'error condicionada pel dispositiu
p_error_cond <- prop.table(xt, margin = 1)[, "error"]

p_error
p_error_cond

# Si són molt diferents, el dispositiu aporta informació
p_error_cond - p_error
\end{lstlisting}
\vspace{0.15cm}
\begin{block}{Important}
Amb una mostra petita, diferències petites poden ser només soroll. La inferència estadística ajudarà a decidir si la diferència és compatible amb l'atzar.
\end{block}
\end{frame}

% ------------------------------------------------------------------
\begin{frame}{El problema dels aniversaris}
\begin{columns}[T]
\column{0.56\textwidth}
Volem calcular la probabilitat que, entre $k$ persones, almenys dues comparteixin aniversari.
\vspace{0.15cm}
\begin{itemize}
  \item És més fàcil calcular el complementari: cap aniversari repetit.
\end{itemize}
\[
\Prob(\text{cap repetit})=\frac{365\cdot 364\cdot \ldots \cdot (365-k+1)}{365^k}.
\]
\[
\Prob(\text{almenys un repetit})=1-\Prob(\text{cap repetit}).
\]
\column{0.38\textwidth}
\begin{center}
\small
\begin{tabular}{cc}
\toprule
$k$ & $\Prob(\text{repetit})$ \\
\midrule
10 & 0.1169 \\
20 & 0.4114 \\
23 & 0.5073 \\
30 & 0.7063 \\
\bottomrule
\end{tabular}
\end{center}
\begin{alertblock}{Sorprenent?}
Amb 23 persones ja superem el 50\%.
\end{alertblock}
\end{columns}
\end{frame}

% ------------------------------------------------------------------
\begin{frame}[fragile]{Simular el problema dels aniversaris}
\begin{lstlisting}
simula_aniversaris <- function(k) {
  dies <- sample(1:365, size = k, replace = TRUE)
  any(duplicated(dies))
}

set.seed(123)
B <- 10000
ks <- c(10, 20, 23, 30)

prob_sim <- sapply(ks, function(k) {
  mean(replicate(B, simula_aniversaris(k)))
})

tibble::tibble(k = ks, prob_simulada = prob_sim)
\end{lstlisting}
\vspace{0.2cm}
\begin{block}{Lectura computacional}
Aquest exemple mostra com una probabilitat aparentment complicada es pot entendre amb una combinació de complementaris i simulació.
\end{block}
\end{frame}

% ------------------------------------------------------------------
\section{Probabilitat total i Bayes}

\begin{frame}{Particions de l'espai mostral}
\begin{columns}[T]
\column{0.52\textwidth}
Una partició divideix $\Omega$ en esdeveniments disjunts:
\[
\Omega=H_1\cup H_2\cup\cdots\cup H_n,
\qquad H_i\cap H_j=\emptyset.
\]
Qualsevol esdeveniment $A$ es pot descompondre com:
\[
A=(A\cap H_1)\cup\cdots\cup(A\cap H_n).
\]
\column{0.42\textwidth}
\begin{center}
\begin{tikzpicture}[scale=0.9]
\draw[thick,rounded corners] (0,0) rectangle (4,2.7);
\fill[UABLight] (0,0) rectangle (1.3,2.7);
\fill[SoftGray] (1.3,0) rectangle (2.7,2.7);
\fill[AccentOrange!18] (2.7,0) rectangle (4,2.7);
\draw[thick] (1.3,0) -- (1.3,2.7);
\draw[thick] (2.7,0) -- (2.7,2.7);
\node at (0.65,2.35) {$H_1$};
\node at (2.0,2.35) {$H_2$};
\node at (3.35,2.35) {$H_3$};
\draw[DeepBlue,very thick] (2,1.35) ellipse (1.45 and 0.7);
\node[DeepBlue] at (2,1.35) {$A$};
\node at (2,-0.35) {$\Omega$};
\end{tikzpicture}
\end{center}
\end{columns}
\end{frame}

% ------------------------------------------------------------------
\begin{frame}{Fórmula de probabilitat total}
Si $H_1,\ldots,H_n$ formen una partició:
\[
\Prob(A)=\Prob(A\cap H_1)+\cdots+\Prob(A\cap H_n).
\]
Com que
\[
\Prob(A\cap H_i)=\Prob(A\mid H_i)\Prob(H_i),
\]
tenim:
\[
\boxed{\Prob(A)=\sum_{i=1}^n \Prob(A\mid H_i)\Prob(H_i)}
\]
\vspace{0.1cm}
\begin{block}{Interpretació}
La probabilitat global de $A$ és una mitjana ponderada de les probabilitats condicionades de $A$ dins de cada grup.
\end{block}
\end{frame}

% ------------------------------------------------------------------
\begin{frame}{Teorema de Bayes}
\begin{block}{Fórmula}
Si $H_1,\ldots,H_n$ formen una partició i observem $A$:
\[
\boxed{
\Prob(H_j\mid A)=
\frac{\Prob(A\mid H_j)\Prob(H_j)}
{\sum_{i=1}^n \Prob(A\mid H_i)\Prob(H_i)}
}
\]
\end{block}
\vspace{0.2cm}
\begin{columns}[T]
\column{0.48\textwidth}
\begin{itemize}
  \item $\Prob(H_j)$: probabilitat abans d'observar $A$.
  \item $\Prob(A\mid H_j)$: versemblança de l'evidència.
  \item $\Prob(H_j\mid A)$: probabilitat després d'observar $A$.
\end{itemize}
\column{0.48\textwidth}
\begin{alertblock}{Lectura informàtica}
Bayes és el nucli de molts raonaments de diagnòstic: spam, alertes, tests automàtics, detecció d'intrusions i classificadors probabilístics.
\end{alertblock}
\end{columns}
\end{frame}

% ------------------------------------------------------------------
\begin{frame}{Exemple: filtre de correu brossa}
\begin{columns}[T]
\column{0.52\textwidth}
Un compte de correu rep missatges:
\begin{itemize}
  \item $30\%$ són spam: $\Prob(S)=0.3$.
  \item Si és spam, el $30\%$ conté una paraula sospitosa $V$: $\Prob(V\mid S)=0.3$.
  \item Si no és spam, l'1\% conté $V$: $\Prob(V\mid \overline{S})=0.01$.
\end{itemize}
Si un missatge conté $V$, quina és la probabilitat que sigui spam?
\column{0.42\textwidth}
\begin{align*}
\Prob(V) &= 0.3\cdot 0.3 + 0.01\cdot 0.7 \\
&=0.097,\\[0.2cm]
\Prob(S\mid V) &= \frac{\Prob(V\mid S)\Prob(S)}{\Prob(V)}\\
&=\frac{0.3\cdot 0.3}{0.097}=0.9278.
\end{align*}
\begin{block}{Conclusió}
La paraula $V$ és molt informativa, però no perfecta.
\end{block}
\end{columns}
\end{frame}

% ------------------------------------------------------------------
\begin{frame}[fragile]{Bayes amb una taula en R}
\begin{lstlisting}
prior_spam <- 0.30
p_v_spam <- 0.30
p_v_no_spam <- 0.01

p_v <- p_v_spam * prior_spam +
       p_v_no_spam * (1 - prior_spam)

posterior_spam <- p_v_spam * prior_spam / p_v
posterior_spam
\end{lstlisting}
\vspace{0.2cm}
\begin{block}{Extensió natural}
En lloc d'una sola paraula, podem tenir moltes característiques: paraules, domini, adjunts, hora d'enviament, historial del remitent... Aquesta és la porta d'entrada als classificadors probabilístics.
\end{block}
\end{frame}

% ------------------------------------------------------------------
\begin{frame}{Exemple: alarma d'un sistema}
\begin{columns}[T]
\column{0.52\textwidth}
Un sistema de monitoratge dispara una alarma si detecta una incidència.
\begin{itemize}
  \item Probabilitat real d'incidència: $\Prob(I)=0.01$.
  \item Si hi ha incidència, l'alarma es dispara amb probabilitat $0.97$.
  \item Si no hi ha incidència, hi ha falsa alarma amb probabilitat $0.02$.
\end{itemize}
Si sona l'alarma, quina és la probabilitat que hi hagi incidència?
\column{0.42\textwidth}
\begin{align*}
\Prob(I\mid A)&=\frac{\Prob(A\mid I)\Prob(I)}{\Prob(A\mid I)\Prob(I)+\Prob(A\mid \overline{I})\Prob(\overline{I})}\\[0.2cm]
&=\frac{0.97\cdot 0.01}{0.97\cdot 0.01+0.02\cdot 0.99}\\[0.2cm]
&=0.3288.
\end{align*}
\end{columns}
\begin{alertblock}{Lliçó}
Una alarma molt bona pot tenir una probabilitat posterior moderada si l'esdeveniment real és molt rar.
\end{alertblock}
\end{frame}

% ------------------------------------------------------------------
\begin{frame}{Base rate fallacy: el risc d'oblidar la prevalença}
\begin{columns}[T]
\column{0.55\textwidth}
\begin{itemize}
  \item L'alarma té sensibilitat alta: $\Prob(A\mid I)=0.97$.
  \item Però la incidència és molt rara: $\Prob(I)=0.01$.
  \item Moltes alarmes provenen de casos sense incidència, perquè són molt més nombrosos.
\end{itemize}
\vspace{0.1cm}
\begin{block}{Traducció informàtica}
En detecció d'intrusions, frau o anomalies, no podem mirar només la taxa d'encert. Cal tenir en compte la freqüència real de l'esdeveniment.
\end{block}
\column{0.4\textwidth}
\centering
\small
\begin{tabular}{lcc}
\toprule
 & Alarma & No alarma \\
\midrule
Incidència & 0.0097 & 0.0003 \\
No incidència & 0.0198 & 0.9702 \\
\midrule
Total & 0.0295 & 0.9705 \\
\bottomrule
\end{tabular}
\vspace{0.2cm}
\[
\frac{0.0097}{0.0295}=0.3288
\]
\end{columns}
\end{frame}

% ------------------------------------------------------------------
\begin{frame}{Exemple: qualitat de software o hardware}
\begin{columns}[T]
\column{0.55\textwidth}
Tres proveïdors fabriquen components. Les quotes i percentatges defectuosos són:
\begin{center}
\small
\begin{tabular}{lcc}
\toprule
Proveïdor & Quota & Defectuosos \\
\midrule
A & 70\% & 4\% \\
B & 24\% & 3\% \\
C & 6\% & 1\% \\
\bottomrule
\end{tabular}
\end{center}
\begin{itemize}
  \item Probabilitat que un component sigui defectuós?
  \item Si és defectuós, probabilitat que sigui de C?
  \item Si funciona, probabilitat que sigui de B?
\end{itemize}
\column{0.4\textwidth}
\begin{align*}
\Prob(D) &= 0.04\cdot 0.70 + 0.03\cdot 0.24 + 0.01\cdot 0.06 \\
&=0.0358,\\[0.2cm]
\Prob(C\mid D) &= \frac{0.01\cdot 0.06}{0.0358}=0.0168,\\[0.2cm]
\Prob(B\mid \overline{D}) &= \frac{0.97\cdot 0.24}{1-0.0358}=0.2414.
\end{align*}
\end{columns}
\end{frame}

% ------------------------------------------------------------------
\begin{frame}[fragile]{Codi genèric per a probabilitat total i Bayes}
\begin{lstlisting}
proveidors <- tibble::tribble(
  ~proveidor, ~quota, ~prob_defecte_cond,
  "A",        0.70,    0.04,
  "B",        0.24,    0.03,
  "C",        0.06,    0.01
)

prob_defecte_total <- sum(
  proveidors$quota * proveidors$prob_defecte_cond
)

proveidors <- proveidors |>
  dplyr::mutate(
    prob_proveidor_si_defecte =
      quota * prob_defecte_cond / prob_defecte_total
  )

proveidors
\end{lstlisting}
\vspace{0.1cm}
\begin{block}{Programar estadística també és programar bé}
Noms clars, funcions petites i informes reproduïbles redueixen errors matemàtics i errors de codi.
\end{block}
\end{frame}

% ------------------------------------------------------------------
\section{Simulació i paral·lelització}

\begin{frame}{Quan la probabilitat exacta és difícil}
\begin{columns}[T]
\column{0.55\textwidth}
En problemes reals, el càlcul exacte pot ser difícil:
\begin{itemize}
  \item molts components,
  \item dependències entre fallades,
  \item regles complexes,
  \item grans volums de dades,
  \item probabilitats rares.
\end{itemize}
\vspace{0.1cm}
\begin{block}{Estratègia}
Simular moltes vegades el sistema i estimar la probabilitat com a freqüència relativa.
\end{block}
\column{0.4\textwidth}
\begin{tikzpicture}[node distance=0.8cm, every node/.style={font=\small}]
\node[draw,rounded corners,fill=UABLight,text width=3.9cm,align=center] (model) {Model del sistema};
\node[draw,rounded corners,fill=SoftGray,below=of model,text width=3.9cm,align=center] (sim) {Simular $B$ vegades};
\node[draw,rounded corners,fill=SoftGray,below=of sim,text width=3.9cm,align=center] (freq) {Comptar esdeveniments};
\node[draw,rounded corners,fill=AccentOrange!15,below=of freq,text width=3.9cm,align=center] (prob) {Estimar probabilitat};
\draw[->,thick,UABGreen] (model) -- (sim);
\draw[->,thick,UABGreen] (sim) -- (freq);
\draw[->,thick,UABGreen] (freq) -- (prob);
\end{tikzpicture}
\end{columns}
\end{frame}

% ------------------------------------------------------------------
\begin{frame}[fragile]{Exemple: fiabilitat d'un sistema amb rèpliques}
Un servei funciona si almenys 2 de 3 servidors estan actius. Cada servidor falla amb probabilitat 0.05.
\begin{lstlisting}
set.seed(123)
B <- 100000

simula_sistema <- function() {
  actiu <- rbinom(n = 3, size = 1, prob = 0.95)
  sum(actiu) >= 2
}

funciona <- replicate(B, simula_sistema())
mean(funciona)
\end{lstlisting}
\vspace{0.1cm}
\begin{block}{Lectura}
Aquest tipus de simulació connecta probabilitat, arquitectura de sistemes i presa de decisions sobre redundància.
\end{block}
\end{frame}

% ------------------------------------------------------------------
\begin{frame}[fragile]{Accelerar simulacions amb paral·lelització}
\begin{columns}[T]
\column{0.58\textwidth}
\begin{lstlisting}
library(future)
library(furrr)

plan(multisession, workers = 4)

B <- 10000
resultats <- future_map_dbl(1:B, function(i) {
  actiu <- rbinom(n = 3, size = 1, prob = 0.95)
  as.numeric(sum(actiu) >= 2)
})

mean(resultats)
\end{lstlisting}
\column{0.38\textwidth}
\begin{block}{Missatge}
La probabilitat proporciona el model; la computació permet repetir-lo moltes vegades.
\end{block}
\begin{alertblock}{Important}
La paral·lelització és útil quan cada simulació és prou costosa. Per problemes petits, el cost de coordinar processos pot no compensar.
\end{alertblock}
\end{columns}
\end{frame}

\end{document}
